\section{Conclusiones}
\label{sec:conclusiones}

En el parcial nos propusimos implementar el PGM-index estático y compararlo con un B$^+$ tree y
con búsqueda binaria. Implementamos el PLA óptimo, el índice recursivo y los dos B$^+$, y los
medimos sobre 14 conjuntos de datos. Sobre la compacidad (P1), la cota $m\le n/(2\varepsilon)$ se
cumplió siempre. Con datos regulares el índice ocupa unos pocos KB para diez millones de llaves;
solo los casos patológicos de cola pesada llegan a la mitad de la cota. Sobre $\varepsilon$ (P2),
el espacio baja varios órdenes de magnitud al aumentarlo, pero el tiempo es mínimo entre
$\varepsilon=16$ y 64. Con valores mayores domina la búsqueda final, y con valores menores el
índice crece sin que la consulta mejore.

Sobre la ventaja práctica (P3), con arreglos grandes ($n=10^7$) el PGM responde unas dos veces
más rápido que el B$^+$ y que la búsqueda binaria, y ocupa cientos de veces menos que el B$^+$ en
los datos regulares. La ventaja se reduce cuando las llaves tienen saltos de cola pesada, porque
$m$ se acerca a la cota, y cuando el arreglo es chico ($n=10^5$). Además de lo
propuesto, implementamos dos versiones dinámicas del índice. La de buffer busca algo más lento
que el estático y conviene solo si las escrituras son raras. La logarítmica inserta a la par del
B$^+$ dinámico, a cambio de búsquedas varias veces más lentas que las del estático.
